
The organization of the simple executive processor {\bf {\em toyRISC}} (described in {\tt toyRISC.v} module of the design listed in Figure \ref{cod2}) is given in Figure \ref{fig43}, where two simple automata, the {\bf RALU} ({\bf R}egisters with {\bf A}rithmetic-{\bf L}ogic {\bf U}nit) subsystem  and  the {\bf Control} subsystem  are loop connected thus closing the 3rd loop which allow the emergence of the processing function. The only complex subsystems are two small modules: the decode circuit (described in {\tt Decode.v} and the interrupt automaton, {\tt intSect.v}.

\section{Code}

The design has a multi-level organization

\subsection{First Level}

The top level

\placedrawing[H]{figures/C1.lp}{{\tt 01\_theSystem}}{cod1}

\subsection{Second Level}

\placedrawing[H]{figures/C2.lp}{{\tt 02\_toyRISC}}{cod2}

\placedrawing[H]{figures/C3.lp}{{\tt 02\_memories}}{cod3}

\subsection{Third Level}

\placedrawing[H]{figures/C5.lp}{{\tt }}{cod5}

\placedrawing[H]{figures/C6.lp}{{\tt }}{cod6}

\placedrawing[H]{figures/C4.lp}{{\tt }}{cod4}

\placedrawing[H]{figures/C7.lp}{{\tt }}{cod7}

\placedrawing[H]{figures/C8.lp}{{\tt }}{cod8}

\subsection{Fourth Level}

\placedrawing[H]{figures/C12.lp}{{\tt }}{cod12}

\placedrawing[H]{figures/C9.lp}{{\tt }}{cod9}

\placedrawing[H]{figures/C10.lp}{{\tt }}{cod10}

\placedrawing[H]{figures/C11.lp}{{\tt }}{cod11}

\section{Assembly}\label{assembly}

{\footnotesize
\begin{shaded*}
\begin{lstlisting}[language=Verilog]
/************************************************************************
File name: 00_toyRISCcodeGenerator.v
Description:
************************************************************************/
    reg [5:0]   RISCopCode          ;
    reg [4:0]   RISCdest            ;
    reg [4:0]   RISCleft            ;
    reg [4:0]   RISCright           ;
    reg [10:0]  RISCvalue           ;
    reg [6:0]   RISCaddrCounter     ;
    reg [6:0]   RISClabelTab[0:128] ;

    `include "0_ISAdefine.vh"

    task hendLine;
        begin
            dut.MEMORIES.PROG_MEM.progMem[RISCaddrCounter] =
                {   RISCopCode  ,
                    RISCdest    ,
                    RISCleft    ,
                    RISCright   ,
                    RISCvalue   }                   ;
            RISCaddrCounter = RISCaddrCounter + 1   ;
        end
    endtask

    // sets RISClabelTab in the first pass associating
    // 'RISCaddrCounter' with 'labelIndex'
    task LB ;
        input [4:0] labelIndex;

        RISClabelTab[labelIndex] = RISCaddrCounter;
    endtask
    // uses the content of RISClabelTab in the second pass for rjump
    task ULB;
        input [4:0] labelIndex;

        {RISCright, RISCvalue}  =
            RISClabelTab[labelIndex] - RISCaddrCounter;
    endtask
    // uses the content of RISClabelTab in the second pass for call
    task CLB;
        input [4:0] labelIndex;

        {RISCright, RISCvalue}  =
            RISClabelTab[labelIndex];
    endtask

    task HALT;
        begin   RISCopCode  = `rjmp;
                RISCdest    = 0     ;
                RISCleft    = 0     ;
                RISCright   = 0     ;
                RISCvalue   = 0     ;
                hendLine            ;
        end
    endtask

    task NOP;
        begin   RISCopCode  = `nop  ;
                RISCdest    = 0     ;
                RISCleft    = 0     ;
                RISCright   = 0     ;
                RISCvalue   = 0     ;
                hendLine            ;
        end
    endtask

    task RJMP;
        input [5:0] label   ;

        begin   RISCopCode  = `rjmp;
                RISCdest    = 0     ;
                RISCleft    = 0     ;
                ULB(label)          ;
                hendLine            ;
        end
    endtask

    task ZJMP;
        input [4:0] leftAddr;
        input [5:0] label   ;

        begin   RISCopCode  = `zjmp     ;
                RISCdest    = 0         ;
                RISCleft    = leftAddr  ;
                ULB(label)              ;
                hendLine                ;
        end
    endtask

    task NZJMP;
        input [4:0] leftAddr;
        input [5:0] label   ;

        begin   RISCopCode  = `nzjmp    ;
                RISCdest    = 0         ;
                RISCleft    = leftAddr  ;
                ULB(label)              ;
                hendLine                ;
        end
    endtask

    task RET;
        input [4:0] leftAddr;

        begin   RISCopCode  = `ret      ;
                RISCdest    = 0         ;
                RISCleft    = 5'b11111  ;
                RISCright   = 0         ;
                RISCvalue   = 0         ;
                hendLine                ;
        end
    endtask

    task AJMP;
        input [10:0] value  ;

        begin   RISCopCode              = `ajmp;
                RISCdest                = 0     ;
                RISCleft                = 0     ;
                {RISCright, RISCvalue}  = value ;
                hendLine                        ;
        end
    endtask

    task CALL;
        input [5:0] label   ;

        begin   RISCopCode  = `call;
                RISCdest    = 0     ;
                RISCleft    = 0     ;
                CLB(label)          ;
                hendLine            ;
        end
    endtask

    task BIGV;
        input [4:0] destAddr    ;
        input [4:0] leftAddr    ;
        input [4:0] rightAddr   ;

        begin   RISCopCode  = `bigv     ;
                RISCdest    = destAddr  ;
                RISCleft    = leftAddr  ;
                RISCright   = rightAddr ;
                RISCvalue   = 11'b0     ;
                hendLine                ;
        end
    endtask

    task ADD;
        input [4:0] destAddr    ;
        input [4:0] leftAddr    ;
        input [4:0] rightAddr   ;

        begin   RISCopCode  = `add      ;
                RISCdest    = destAddr  ;
                RISCleft    = leftAddr  ;
                RISCright   = rightAddr ;
                RISCvalue   = 11'b0     ;
                hendLine                ;
        end
    endtask

    task SUB;
        input [4:0] destAddr    ;
        input [4:0] leftAddr    ;
        input [4:0] rightAddr   ;

        begin   RISCopCode  = `sub      ;
                RISCdest    = destAddr  ;
                RISCleft    = leftAddr  ;
                RISCright   = rightAddr ;
                RISCvalue   = 11'b0     ;
                hendLine                ;
        end
    endtask

    task CRADD;
        input [4:0] destAddr    ;
        input [4:0] leftAddr    ;
        input [4:0] rightAddr   ;

        begin   RISCopCode  = `addcr    ;
                RISCdest    = destAddr  ;
                RISCleft    = leftAddr  ;
                RISCright   = rightAddr ;
                RISCvalue   = 11'b0     ;
                hendLine                ;
        end
    endtask

    task CRSUB;
        input [4:0] destAddr    ;
        input [4:0] leftAddr    ;
        input [4:0] rightAddr   ;

        begin   RISCopCode  = `subcr    ;
                RISCdest    = destAddr  ;
                RISCleft    = leftAddr  ;
                RISCright   = rightAddr ;
                RISCvalue   = 11'b0     ;
                hendLine                ;
        end
    endtask

    task LEFTSH; // logic left shift
        input [4:0] destAddr    ;
        input [4:0] leftAddr    ;

        begin   RISCopCode  = `lsh      ;
                RISCdest    = destAddr  ;
                RISCleft    = leftAddr  ;
                RISCright   = 5'b0      ;
                RISCvalue   = 11'b0     ;
                hendLine                ;
        end
    endtask

    task ARITHSH; // arithmetic shift
        input [4:0] destAddr;
        input [4:0] leftAddr;

        begin   RISCopCode  = `ash      ;
                RISCdest    = destAddr  ;
                RISCleft    = leftAddr  ;
                RISCright   = 0         ;
                RISCvalue   = 0         ;
                hendLine                ;
        end
    endtask

    task MOVE;
        input [4:0] destAddr;
        input [4:0] leftAddr;

        begin   RISCopCode  = `move ;
                RISCdest    = destAddr  ;
                RISCleft    = leftAddr  ;
                RISCright   = 0         ;
                RISCvalue   = 0         ;
                hendLine                ;
        end
    endtask

    task MULT;
        input [4:0] destAddr    ;
        input [4:0] leftAddr    ;
        input [4:0] rightAddr   ;

        begin   RISCopCode  = `mult ;
                RISCdest    = destAddr  ;
                RISCleft    = leftAddr  ;
                RISCright   = rightAddr ;
                RISCvalue   = 0         ;
                hendLine                ;
        end
    endtask

    task AND;
        input [4:0] destAddr    ;
        input [4:0] leftAddr    ;
        input [4:0] rightAddr   ;

        begin   RISCopCode  = `bwand    ;
                RISCdest    = destAddr  ;
                RISCleft    = leftAddr  ;
                RISCright   = rightAddr ;
                RISCvalue   = 0         ;
                hendLine                ;
        end
    endtask

    task OR;
        input [4:0] destAddr    ;
        input [4:0] leftAddr    ;
        input [4:0] rightAddr   ;

        begin   RISCopCode  = `bwor ;
                RISCdest    = destAddr  ;
                RISCleft    = leftAddr  ;
                RISCright   = rightAddr ;
                RISCvalue   = 0         ;
                hendLine                ;
        end
    endtask

    task XOR;
        input [4:0] destAddr    ;
        input [4:0] leftAddr    ;
        input [4:0] rightAddr   ;

        begin   RISCopCode  = `bwxor    ;
                RISCdest    = destAddr  ;
                RISCleft    = leftAddr  ;
                RISCright   = rightAddr ;
                RISCvalue   = 11'b0     ;
                hendLine                ;
        end
    endtask

    task READ;
        input [4:0] destAddr;
        input [4:0] rightAddr;

        begin   RISCopCode  = `read ;
                RISCdest    = destAddr  ;
                RISCleft    = 5'b0      ;
                RISCright   = rightAddr ;
                RISCvalue   = 11'b0     ;
                hendLine                ;
        end
    endtask

    task STORE;
        input [4:0] leftAddr    ;
        input [4:0] rightAddr   ;

        begin   RISCopCode  = `store    ;
                RISCdest    = 5'b0      ;
                RISCleft    = leftAddr  ;
                RISCright   = rightAddr ;
                RISCvalue   = 11'b0     ;
                hendLine                ;
        end
    endtask

    task VAL;
        input [4:0]     destAddr;
        input [15:0]    value   ;

        begin   RISCopCode  = `val              ;
                RISCdest    = destAddr          ;
                RISCleft    = 5'b0              ;
                {RISCright, RISCvalue} = value  ;
                hendLine                        ;
        end
    endtask

    task EINT; // enable interrupt
        begin   RISCopCode  = `ei   ;
                RISCdest    = 0     ;
                RISCleft    = 0     ;
                RISCright   = 0     ;
                RISCvalue   = 0     ;
                hendLine            ;
        end
    endtask

    task DINT; // disable interrupt
        begin   RISCopCode  = `di   ;
                RISCdest    = 0     ;
                RISCleft    = 0     ;
                RISCright   = 0     ;
                RISCvalue   = 0     ;
                hendLine            ;
        end
    endtask

    // RUNNING
    initial begin   RISCaddrCounter = 0;
                    `include "00_toyRISCprogram.v" // first pass
                    RISCaddrCounter = 0;
                    `include "00_toyRISCprogram.v" // second pass
            end
 \end{lstlisting}
\end{shaded*}
}


\section{Simulation}

%\pagebreak

{\footnotesize
\begin{shaded*}
\begin{lstlisting}[language=Verilog]
/************************************************************************
File name: 00_toyRiscSimulator.v
Description:
************************************************************************/
//`include "DEFINES.vh"
module toyRiscSimulator;
    reg         reset, interrupt, clk   ;
    reg [15:0]  intAddr                 ;
    wire        inta                    ;

    integer i;

    initial begin               clk = 0     ;
                    forever #1  clk = ~clk  ;
            end

    `include "00_toyRISCcodeGenerator.v"

    initial for (i=0; i<32; i=i+1)
        $display("programMemory[%0d] \t = %b",
                    i, dut.MEMORIES.PROG_MEM.progMem[i]);

    initial
        begin       reset       = 1         ;
                    intAddr     = 16'b1111  ;
                    interrupt   = 0         ;
            #4      reset       = 0         ;
            #20     interrupt   = 1         ;
            #150    begin
                    // DISPLAY THE CONTENT OF THE REGISTER FILE
                        for (i=0; i<8; i=i+1)
                            $display("RF[%0d] \t = %b",
                            i, dut.PROCESSOR.RALU.regFile.registers[i]);
                    // DISPLAY THE DATA MEMORY OF HOST
                        for (i=0; i<8; i=i+1)
                            $display("dataMemory[%0d] \t = %b",
                            i, dut.MEMORIES.DATA_MEM.mem[i]);
                    end
            #2      $stop   ;
        end

    theSystem dut(reset, interrupt, intAddr, inta, clk);

    initial begin
        $monitor
        ("t=%0d \t pc=%0d \t inta=%0d \t
        RF[0]=%0d \t RF[1]=%0d\t RF[2]=%0d \t RF[3]=%0d  ...
        \t RF[31]=%0d",
         $time,
         dut.PROCESSOR.Control.pc,
         dut.PROCESSOR.Control.address,
         dut.PROCESSOR.inta,
         dut.PROCESSOR.RALU.regFileIn,
         dut.PROCESSOR.RALU.regFile.registers[0],
         dut.PROCESSOR.RALU.regFile.registers[1],
         dut.PROCESSOR.RALU.regFile.registers[2],
         dut.PROCESSOR.RALU.regFile.registers[3],
         dut.PROCESSOR.RALU.regFile.registers[31]);
    end
endmodule
\end{lstlisting}
\end{shaded*}
}

\placedrawing[H]{figures/C13.lp}{{\tt }}{cod13}

{\footnotesize
\begin{verbatim}
***************************************************************************************
    The binary form of program provided by 00_toyRISCcodeGenerator.v
***************************************************************************************
programMemory[0]     = 01011100000000000000000000001011
programMemory[1]     = 01011100001000000000000000000001
programMemory[2]     = 01010100000000000000000000000000
programMemory[3]     = 11001100000000000000100000000000
programMemory[4]     = 00001100000000001111111111111111
programMemory[5]     = 00011100000000000000000000001010
programMemory[6]     = 00000000000000000000000000000000
programMemory[7]     = 00000000000000000000000000000000
programMemory[8]     = 01011100010000000000000000100001
programMemory[9]     = 00000100000000000000000000000000
programMemory[10]    = 01011100000000000000000001000010
programMemory[11]    = 00010100000111110000000000000000
programMemory[12]    = 00000000000000000000000000000000
programMemory[13]    = 00000000000000000000000000000000
programMemory[14]    = 00000000000000000000000000000000
programMemory[15]    = 00000000000000000000000000000000
programMemory[16]    = 01011100011000000000000000001101
programMemory[17]    = 00000000000000000000000000000000
programMemory[18]    = 00010100000111110000000000000000

***************************************************************************************
    The result of simulation provided by the monitor run by 00_toyRiscSimulator.v
***************************************************************************************
t=0     pc=x     inta=0   RF[0]=x    RF[1]=x   RF[2]=x   RF[3]=x  ...   RF[31]=x
t=1     pc=65535 inta=0   RF[0]=x    RF[1]=x   RF[2]=x   RF[3]=x  ...   RF[31]=x
t=5     pc=0     inta=0   RF[0]=11   RF[1]=x   RF[2]=x   RF[3]=x  ...   RF[31]=x
t=7     pc=1     inta=0   RF[0]=11   RF[1]=x   RF[2]=x   RF[3]=x  ...   RF[31]=x
t=9     pc=2     inta=0   RF[0]=11   RF[1]=1   RF[2]=x   RF[3]=x  ...   RF[31]=x
t=11    pc=3     inta=0   RF[0]=11   RF[1]=1   RF[2]=x   RF[3]=x  ...   RF[31]=x
t=13    pc=4     inta=0   RF[0]=10   RF[1]=1   RF[2]=x   RF[3]=x  ...   RF[31]=x
t=15    pc=3     inta=0   RF[0]=10   RF[1]=1   RF[2]=x   RF[3]=x  ...   RF[31]=x
t=17    pc=4     inta=0   RF[0]=9    RF[1]=1   RF[2]=x   RF[3]=x  ...   RF[31]=x
t=19    pc=3     inta=0   RF[0]=9    RF[1]=1   RF[2]=x   RF[3]=x  ...   RF[31]=x
t=21    pc=4     inta=0   RF[0]=8    RF[1]=1   RF[2]=x   RF[3]=x  ...   RF[31]=x
t=23    pc=3     inta=0   RF[0]=8    RF[1]=1   RF[2]=x   RF[3]=x  ...   RF[31]=x
t=24    pc=3     inta=1   RF[0]=8    RF[1]=1   RF[2]=x   RF[3]=x  ...   RF[31]=x
t=25    pc=15    inta=0   RF[0]=8    RF[1]=1   RF[2]=x   RF[3]=x  ...   RF[31]=4
t=27    pc=16    inta=0   RF[0]=8    RF[1]=1   RF[2]=x   RF[3]=x  ...   RF[31]=4
t=29    pc=17    inta=0   RF[0]=8    RF[1]=1   RF[2]=x   RF[3]=13  ...  RF[31]=4
t=31    pc=18    inta=0   RF[0]=8    RF[1]=1   RF[2]=x   RF[3]=13  ...  RF[31]=4
t=33    pc=4     inta=0   RF[0]=8    RF[1]=1   RF[2]=x   RF[3]=13  ...  RF[31]=4
t=35    pc=3     inta=0   RF[0]=8    RF[1]=1   RF[2]=x   RF[3]=13  ...  RF[31]=4
t=37    pc=4     inta=0   RF[0]=7    RF[1]=1   RF[2]=x   RF[3]=13  ...  RF[31]=4
t=39    pc=3     inta=0   RF[0]=7    RF[1]=1   RF[2]=x   RF[3]=13  ...  RF[31]=4
t=41    pc=4     inta=0   RF[0]=6    RF[1]=1   RF[2]=x   RF[3]=13  ...  RF[31]=4
t=43    pc=3     inta=0   RF[0]=6    RF[1]=1   RF[2]=x   RF[3]=13  ...  RF[31]=4
t=45    pc=4     inta=0   RF[0]=5    RF[1]=1   RF[2]=x   RF[3]=13  ...  RF[31]=4
t=47    pc=3     inta=0   RF[0]=5    RF[1]=1   RF[2]=x   RF[3]=13  ...  RF[31]=4
t=49    pc=4     inta=0   RF[0]=4    RF[1]=1   RF[2]=x   RF[3]=13  ...  RF[31]=4
t=51    pc=3     inta=0   RF[0]=4    RF[1]=1   RF[2]=x   RF[3]=13  ...  RF[31]=4
t=53    pc=4     inta=0   RF[0]=3    RF[1]=1   RF[2]=x   RF[3]=13  ...  RF[31]=4
t=55    pc=3     inta=0   RF[0]=3    RF[1]=1   RF[2]=x   RF[3]=13  ...  RF[31]=4
t=57    pc=4     inta=0   RF[0]=2    RF[1]=1   RF[2]=x   RF[3]=13  ...  RF[31]=4
t=59    pc=3     inta=0   RF[0]=2    RF[1]=1   RF[2]=x   RF[3]=13  ...  RF[31]=4
t=61    pc=4     inta=0   RF[0]=1    RF[1]=1   RF[2]=x   RF[3]=13  ...  RF[31]=4
t=63    pc=3     inta=0   RF[0]=1    RF[1]=1   RF[2]=x   RF[3]=13  ...  RF[31]=4
t=65    pc=4     inta=0   RF[0]=0    RF[1]=1   RF[2]=x   RF[3]=13  ...  RF[31]=4
t=67    pc=5     inta=0   RF[0]=0    RF[1]=1   RF[2]=x   RF[3]=13  ...  RF[31]=4
t=69    pc=10    inta=0   RF[0]=0    RF[1]=1   RF[2]=x   RF[3]=13  ...  RF[31]=6
t=71    pc=11    inta=0   RF[0]=66   RF[1]=1   RF[2]=x   RF[3]=13  ...  RF[31]=6
t=73    pc=6     inta=0   RF[0]=66   RF[1]=1   RF[2]=x   RF[3]=13  ...  RF[31]=6
t=75    pc=7     inta=0   RF[0]=66   RF[1]=1   RF[2]=x   RF[3]=13  ...  RF[31]=6
t=77    pc=8     inta=0   RF[0]=66   RF[1]=1   RF[2]=x   RF[3]=13  ...  RF[31]=6
t=79    pc=9     inta=0   RF[0]=66   RF[1]=1   RF[2]=33  RF[3]=13  ...  RF[31]=6
\end{verbatim}
}

%\placedrawing[H]{C14.lp}{{\tt }}{cod14}

%\placedrawing[H]{C15.lp}{{\tt }}{cod15} 